\chapter[Background]{Background}
\label{ch:master-background}

\section{Introduction}

This chapter establishes the concepts and review method used to reason about
trustworthy recovery of software components from compiled Android
applications. The thesis treats a recovered \SBOM as a set of evidence-backed
claims rather than as a simple list of names and versions. This distinction
requires a common vocabulary for component identity, provenance, uncertainty,
and transformation risk, followed by a transparent method for selecting and
comparing prior research.

\section{Software Supply Chains and SBOMs}

Modern applications assemble code, resources, native binaries, generated
artifacts, and transitive dependencies from many independently maintained
projects. An \SBOM records these components so that operators can reason about
inventory, licensing, provenance, and known vulnerabilities. Standards define
how such information is represented, but a representation format does not
establish that the recorded component is actually present or that its exact
version is known. This distinction is consistent with supply-chain guidance
that treats component transparency as one control within a broader risk
management process, and with empirical SBOM research showing persistent
generation, consistency, and consumption challenges
\cite{ntia2021minimum,nist2022scrm,xia2023sbom}.

This thesis considers both SPDX and CycloneDX, while using Package URL as a
standard package-identity model
\cite{spdxSpecification,cyclonedxSpecification,packageUrlSpecification}.
These standards therefore define the target vocabulary and exchange
structures, while the literature synthesis defines the evidential conditions
under which individual fields may be populated responsibly.

\subsection{Authoritative and Reconstructed SBOMs}

An authoritative build-time \SBOM can use dependency resolution metadata,
lockfiles, build graphs, and package-manager records. By contrast, an
artifact-reconstructed \SBOM begins after those authoritative records have
been removed or made unavailable. It must infer a component from observable
residue in an \APK, compare that residue with references, and preserve the
uncertainty introduced by compilation and incomplete coverage.

The distinction is central to this thesis. A reconstructed component identity
may be credible even when its exact version is unresolved. Similarly, a
syntactically valid CycloneDX or SPDX document may contain semantically weak
claims. Trustworthiness is consequently a property of each claim and its
evidence trail, not of the file format alone.

\section{Android Applications and Build Cycle}

Android applications are built from source code, framework components,
resources, and packaged dependencies. The platform exposes a clear execution
model to the developer, but that model is partly transformed before the
application reaches an analyst. Understanding the development cycle is
therefore essential before discussing binary recovery.

\subsection{Application Components}

Android applications are organized around a small set of component types:
activities for user-facing screens, services for background work, broadcast
receivers for event-driven actions, and content providers for structured data
sharing \cite{androidAppFundamentals}. These components are declared in the
manifest, linked to permissions and intent filters, and connected to the
Android framework through runtime callbacks.

The manifest is important because it records the application package name,
component declarations, permissions, SDK constraints, and other configuration
data \cite{androidManifestOverview}. However, it does not by itself represent
the full dependency graph of the application. A manifest can therefore support
contextual interpretation, but it cannot be treated as a complete SBOM.

\subsection{Development, Build, Packaging, and Release}

Android development typically uses Gradle and the Android Gradle Plugin to
resolve source modules and external libraries before compilation
\cite{androidBuildOverview,gradleVariantResolution}. A dependency may resolve
to different physical artifacts depending on platform, build type, flavor, or
metadata-only packaging. The build cycle usually proceeds through source
compilation, resource processing, bytecode conversion, packaging, signing,
and distribution.

Java and Kotlin source code is first compiled to JVM bytecode. D8 then
converts that bytecode into DEX form for the Android runtime, while R8 may
shrink, optimize, and obfuscate the release build \cite{androidD8,androidShrinkCode}. The final
artifact therefore reflects both the build configuration and the optimization
policy, not only the original source tree.

In this context, the \emph{source project} denotes the developer-visible
workspace: source files, Gradle scripts, manifests, resources, and declared
dependencies. The \emph{dependency graph} is the resolved set of direct and
transitive libraries selected for a concrete build variant. A build variant is
important because debug and release builds, Android API targets, flavors, and
platform attributes can lead to different resolved artifacts. The build graph
is therefore stronger evidence than the final APK for declaring an SBOM, but
it is not always available to an external analyst.

Resolved libraries usually arrive as JAR or AAR artifacts. A JAR mainly
contains JVM bytecode and metadata, whereas an AAR is Android-specific and may
also contain resources, a manifest fragment, assets, and native libraries.
During packaging, these artifacts are merged with the application code.
Compilation and conversion then change their representation: D8 emits Android
DEX bytecode, and R8 can remove unused classes, inline methods, rename
identifiers, merge classes, and rewrite call sites. \emph{Shrinking} removes
unused code and resources, while \emph{obfuscation} makes names less
recognizable. The signed APK is the distribution artifact after these
transformations and after integrity/authenticity metadata has been added.

Figure~\ref{fig:master-android-build-stages} summarizes this development and
release path. The figure is important because it separates the places where an
ideal build-time SBOM can be produced from the later artifact that binary
recovery must inspect.

\begin{figure}[H]
\centering
\resizebox{0.96\textwidth}{!}{%
\begin{tikzpicture}[
  node distance=7mm and 8mm,
  box/.style={draw, rounded corners, align=center, text width=0.18\textwidth,
    minimum height=11mm, fill=blue!5},
  transform/.style={draw, rounded corners, align=center, text width=0.17\textwidth,
    minimum height=11mm, fill=warningorange!14},
  output/.style={draw, rounded corners, align=center, text width=0.18\textwidth,
    minimum height=11mm, fill=evidencegreen!10},
  arrow/.style={-{Latex[length=2.5mm]}, thick}
]
\node[box] (project) {Source project\\code, manifest, resources};
\node[box, right=of project] (deps) {Resolved dependency graph\\direct and transitive};
\node[transform, right=of deps] (compile) {Compile and convert\\JVM bytecode to DEX};
\node[transform, right=of compile] (opt) {R8 release transforms\\shrink, optimize, obfuscate};
\node[output, right=of opt] (apk) {Signed APK\\distribution artifact};

\draw[arrow] (project) -- (deps);
\draw[arrow] (deps) -- (compile);
\draw[arrow] (compile) -- (opt);
\draw[arrow] (opt) -- (apk);
\end{tikzpicture}
}
\caption{Android development and release stages from source project to signed APK, based on Android build documentation \cite{androidBuildOverview,androidD8,androidShrinkCode}.}
\label{fig:master-android-build-stages}
\end{figure}

\subsection{APK Structure and DEX}

An \APK is a ZIP-format archive that usually contains one or more DEX files,
the manifest, compiled resources, raw assets, native libraries, and signing
material \cite{androidApkAnalyzer}. The DEX files contain types, strings,
fields, methods, classes, annotations, and instructions. When an application
exceeds the single-DEX reference limit, multidex packaging introduces files
such as \texttt{classes.dex} and \texttt{classes2.dex}
\cite{androidMultidex}.

This structure matters for recovery because each region contributes different
evidence. The manifest provides configuration context, DEX files expose code
and structural relationships, resources may reveal package or product hints,
assets may contain bundled metadata, and native libraries may reveal
additional files that are relevant to provenance or vulnerability analysis.
None of these regions is a complete inventory by itself.

Figure~\ref{fig:master-apk-anatomy} shows the APK as a set of complementary
evidence regions. DEX is the main executable evidence because it contains the
classes, methods, strings, fields, and instructions that remain after
compilation. The manifest is declarative evidence: it names components,
permissions, SDK constraints, package identity, and other build-time decisions.
Resources and assets are packaging evidence; they may preserve names,
configuration files, or markers from libraries even when Java/Kotlin names
have been changed. Native files, typically \texttt{.so} libraries, are binary
objects that may have their own provenance and vulnerability surface. Signing
metadata authenticates the APK package but does not prove the correctness of
any recovered component claim.

Two additional concepts are important for structured evidence mapping. First,
\emph{references} denote external method, field, or type calls within the DEX
bytecode that target external packages or standard Android libraries. These
references survive renaming and shrinking because R8 must preserve call site
targets to maintain functionality. Second, \emph{build metadata} includes
declared SDK levels, version codes, build variant configurations, and signing
material. While this metadata is easy to modify, it provides essential context
for verifying the integrity and timeline of the analyzed APK.

\begin{figure}[H]
\centering
\resizebox{0.90\textwidth}{!}{%
\begin{tikzpicture}[
  node distance=6mm and 7mm,
  center/.style={draw, rounded corners, thick, align=center, text width=0.22\textwidth,
    minimum height=12mm, fill=evidencegreen!10},
  part/.style={draw, rounded corners, align=center, text width=0.22\textwidth,
    minimum height=11mm, fill=blue!5},
  arrow/.style={-{Latex[length=2mm]}, thick}
]
\node[center] (apk) {APK archive\\signed distribution package};
\node[part, above left=of apk] (dex) {DEX files\\classes, methods, opcodes};
\node[part, above right=of apk] (manifest) {Manifest\\components, permissions, SDK};
\node[part, below left=of apk] (resources) {Resources and assets\\layouts, strings, metadata};
\node[part, below right=of apk] (native) {Native files\\\texttt{.so} binaries and ABI paths};

\draw[arrow] (apk) -- (dex);
\draw[arrow] (apk) -- (manifest);
\draw[arrow] (apk) -- (resources);
\draw[arrow] (apk) -- (native);
\end{tikzpicture}
}
\caption{Main APK evidence regions used by static SBOM recovery, adapted from APK analysis documentation \cite{androidApkAnalyzer,androidManifestOverview}.}
\label{fig:master-apk-anatomy}
\end{figure}

\subsection{Libraries and Dependencies}

Android applications reuse both direct and transitive dependencies. A direct
dependency is declared explicitly in the build; a transitive dependency is
pulled in by another dependency. Variant-aware resolution can select Android,
JVM, Kotlin, or metadata artifacts for the same coordinate
\cite{gradleVariantResolution}. That is why the binary artifact may differ
from the original source dependency declaration even when the build is
correct.

Libraries are commonly distributed as JAR or AAR files before compilation
\cite{androidLibrary}. Their coordinates are usually expressed in Maven form,
which makes them useful for SBOM generation and vulnerability queries.
However, once those artifacts have been merged, transformed, or shaded into an
APK, the library boundary is no longer directly visible. Binary recovery must
therefore infer whether an observed feature belongs to a library family, a
specific version, or only an uncertain candidate set.

\subsection{Static and Dynamic Analysis}

Static analysis inspects the APK without executing it, using archive parsing,
manifest decoding, resource inspection, DEX disassembly, and structural
feature extraction \cite{owaspStaticAnalysis}. Dynamic analysis observes the
application while it runs, which can reveal runtime-loaded code or behavior
that is invisible in the static archive \cite{owaspReverseEngineering}. Both
approaches are useful, but they answer different questions.

This thesis focuses on static analysis because the goal is to recover
evidence from the delivered artifact itself. Dynamic analysis remains relevant
as a future extension or validation aid, especially for encrypted payloads,
custom class loaders, and other runtime-only behaviors
\cite{androidDynamicCode}.

\section{Trust, Evidence, Provenance, and Uncertainty}

\subsection{Evidence and Claims}

An evidence item is an observable fact used to support or challenge a
component claim. Examples include a package fragment, an opcode fingerprint, a
class-dependency relation, a resource marker, a native-library name, or a
source-validated expected dependency. A claim is the interpretation built
from one or more such items, for example that a library family is present or
that a particular version is the best-supported candidate.

In this context, a \emph{bounded claim} refers to a recovered SBOM assertion
whose certainty and completeness are strictly limited by the available binary
evidence and compiler transformations. Unlike a build-time SBOM, which makes
authoritative assertions based on complete project declarations, a recovered
SBOM can make only partial claims that are bounded by what survived shrinking,
obfuscation, and method inlining.

\subsection{Provenance and Independence}

Provenance records where evidence came from, how it was extracted, and which
transformations may have affected it. Two matching fields are not independent
when they are derived from the same underlying artifact feature. Confidence
must therefore not be increased merely by counting correlated signals as
separate confirmations. The framework adopts the W3C PROV distinction between
entities, activities, and agents as a general provenance vocabulary, while
specializing it for APK observations, reference profiles, validation
artifacts, and analysis steps \cite{moreau2013prov}.

\subsection{Uncertainty and Abstention}

Uncertainty can arise from missing evidence, conflicting evidence, incomplete
reference coverage, adjacent versions with near-identical code, partial
library inclusion, or destructive optimization. The framework allows
ambiguity sets and abstention because an explicit unresolved result is more
auditable than a forced exact version. Reject-option research provides the
general decision-theoretic precedent for withholding a classification when
the expected error cost is too high \cite{chow1970reject}. This policy also limits downstream
vulnerability authority: a vulnerability finding cannot be stronger than the
identity and version evidence on which it depends.

\section{Android Compiled-Artifact Context}

The build-cycle concepts above explain why compiled Android evidence is
structurally weaker than source-level dependency data. Android compilation
and R8 transformations directly affect the stability and interpretability of
recovered evidence \cite{androidShrinkCode}.

Java or Kotlin source is compiled and transformed into \DEX bytecode, while
resources and manifests are packaged into an \APK. Android applications
therefore carry both executable code and declarative metadata: the final
artifact can still expose components, permissions, SDK constraints, resource
names, and native binaries, but those signals no longer appear in the same
shape as the original source tree \cite{androidAppFundamentals,androidBuildOverview,androidManifestOverview,androidApkAnalyzer}. Dependency boundaries that were explicit
during the build may disappear when classes are merged into one or more DEX
files. R8 can rename identifiers, remove unused code, rewrite control flow,
merge classes, inline methods, and change call sites. Packaging may also
introduce multidex layouts, shaded or relocated code, native libraries,
generated classes, and dynamically loaded components.

These transformations do not affect all evidence equally. Names and package
trees are inexpensive but fragile. Opcode, control-flow, call, and dependency
structures can survive renaming, yet they are still altered by shrinking and
optimization. Resources and metadata can provide independent support, but they
may identify only a family rather than an exact version. The state-of-the-art
taxonomy in \cref{ch:master-sota} therefore classifies both evidence sources
and the transformations that threaten them.

Figure~\ref{fig:master-android-build-evidence} illustrates why binary
recovery is weaker than build-time inventory generation. The dependency graph
is visible before compilation, but the audit target is often the transformed
APK after boundaries, names, and unused code have changed.

\begin{figure}[H]
\centering
\resizebox{\textwidth}{!}{%
\begin{tikzpicture}[
  node distance=7mm and 8mm,
  stage/.style={draw, rounded corners, align=center, text width=0.20\textwidth,
    minimum height=11mm, fill=blue!5},
  transform/.style={draw, rounded corners, align=center, text width=0.19\textwidth,
    minimum height=11mm, fill=warningorange!14},
  evidence/.style={draw, rounded corners, align=center, text width=0.22\textwidth,
    minimum height=11mm, fill=evidencegreen!10},
  arrow/.style={-{Latex[length=2.5mm]}, thick}
]
\node[stage] (source) {Source project\\and dependency graph};
\node[stage, right=of source] (resolved) {Resolved JAR/AAR\\artifacts};
\node[transform, right=of resolved] (compile) {D8/R8 conversion,\\shrinking, obfuscation};
\node[stage, right=of compile] (apk) {Signed APK\\DEX, resources, native files};
\node[evidence, below=of source] (buildsbom) {Build-time SBOM\\declared/resolved metadata};
\node[evidence, below=of apk] (recovered) {Recovered SBOM\\observed bounded claims};

\draw[arrow] (source) -- (resolved);
\draw[arrow] (resolved) -- (compile);
\draw[arrow] (compile) -- (apk);
\draw[arrow] (source) -- (buildsbom);
\draw[arrow] (apk) -- (recovered);
\draw[arrow, dashed] (buildsbom) -- node[below, font=\scriptsize, align=center]
  {related but\\not identical} (recovered);
\end{tikzpicture}
}
\caption{Android build and evidence-loss path from resolved dependencies to a recovered SBOM claim, adapted from Android build and APK analysis documentation \cite{androidBuildOverview,androidApkAnalyzer}.}
\label{fig:master-android-build-evidence}
\end{figure}

\section{Software Component Attribution}

\subsection{Identity, Version, and Security Authority}

Component attribution contains at least three distinct decisions:

\begin{enumerate}
  \item \textbf{identity attribution}: whether a component family or
    coordinate is present;
  \item \textbf{version attribution}: which exact version or bounded version
    set is supported; and
  \item \textbf{security interpretation}: which vulnerability statements are
    permitted by the preceding evidence.
\end{enumerate}

A method that performs well at library-family detection does not
automatically provide exact-version authority. Likewise, a vulnerability
database match is not confirmation that the affected version is present.
Keeping these decisions separate prevents a plausible identity match from
being silently promoted into an exact and security-critical statement.

Figure~\ref{fig:master-claim-authority-chain} shows the authority chain. Each
step depends on the previous one: vulnerability interpretation is meaningful
only when component identity and version evidence are strong enough for the
affected-range query being made.

\begin{figure}[H]
\centering
\begin{tikzpicture}[
  node distance=8mm,
  box/.style={draw, rounded corners, align=center, text width=0.78\textwidth,
    minimum height=9mm, fill=blue!5},
  warn/.style={draw, rounded corners, align=center, text width=0.78\textwidth,
    minimum height=9mm, fill=warningorange!13},
  arrow/.style={-{Latex[length=2.5mm]}, thick}
]
\node[box] (evidence) {APK evidence: DEX, manifest, resources, native files, references};
\node[box, below=of evidence] (identity) {Component identity claim};
\node[box, below=of identity] (version) {Exact version or bounded version-set claim};
\node[box, below=of version] (range) {Affected-range query in vulnerability data};
\node[warn, below=of range] (authority) {Permitted security interpretation};

\draw[arrow] (evidence) -- (identity);
\draw[arrow] (identity) -- (version);
\draw[arrow] (version) -- (range);
\draw[arrow] (range) -- (authority);
\end{tikzpicture}
\caption{Authority chain linking APK evidence to vulnerability interpretation.}
\label{fig:master-claim-authority-chain}
\end{figure}

\subsection{Reference Dependence}

Many attribution methods compare an application with profiles extracted from
known library releases. Such a method can recognize only what its reference
corpus represents. Failure may therefore result from the matcher, the
extractor, or a missing reference rather than from absence of the component.
Reference coverage is treated as a first-class validity condition throughout
the thesis.

\section{Evidence Sources in APKs}

The review extracts evidence into recurring families:

\begin{itemize}
  \item lexical and namespace evidence, including package and class names;
  \item bytecode and semantic evidence, including opcodes, API references,
    literals, and method descriptors;
  \item structural evidence, including control-flow, call, inheritance, and
    class-dependency graphs;
  \item resource and packaging evidence, including manifests, resource names,
    assets, and native-library files;
  \item external reference and metadata evidence, including library profiles,
    Maven coordinates, dependency relations, and release chronology;
  \item learned representations derived from bytecode or structural features;
  \item validation evidence, including controlled builds, source manifests,
    and source-based observations.
\end{itemize}

No family is assumed to be universally authoritative. The framework evaluates
its origin, transformation resilience, version specificity, independence,
coverage, and validation status.

\section{Conclusion}

This chapter separated a valid SBOM representation from the evidential
strength of its contents and defined component attribution as a sequence of
identity, version, and security-authority decisions. It also froze a
structured mapping protocol that prevents unavailable full texts,
incompatible metrics, and correlated signals from being silently promoted
into strong claims.

These background concepts explain why Android APK-to-SBOM recovery cannot be
evaluated through a single accuracy value. The same output may be trustworthy
as a library-family hint, weak as an exact-version assertion, and unsafe as a
vulnerability conclusion. A reliable review must therefore preserve the
evidence source, transformation setting, reference coverage, denominator, and
validation context attached to each reported result. The next chapter
summarizes the frozen review protocol and applies these rules to classify the
state of the art and identify the remaining research gaps.
